
Multirotor vehicles have recently gained widespread use in various
application areas, such as transportation, surveillance, and even
entertainment. These platforms are mechanically simple and typically
lightweight, resulting in a low cost of acquisition and
maintenance. However, to the best knowledge of the authors, there is
no known application of multirotor vehicles in space
environments. Discarding the obvious impossibility of maneuvering in
the absence of air, we consider here the use of these vehicles in
human-compatible pressurised spaces under microgravity, such as inside
inhabited space stations.

The idea of aerial vehicles inside space stations is not new. The NASA
project SPHERES (Synchronized Position Hold Engage Reorient
Experimental Satellites) started in 2000 with the design of a small
pressurised air propulsion vehicle~\cite{miller00}. In 2006 three
SPHERES units were deployed aboard the International Space Station
(ISS)~\cite{nolet07}.  However, the use of pressurised air makes the
mechanical complexity significantly higher than a multirotor based
solution.  More recently, NASA proposed the Astrobee vehicle, with a
simpler propulsion system based on several centrifugal fans and
nozzles~\cite{bualat15}. The Astrobee also features a 2 DoF arm and
docking capability. However, most of its volume (about 67\%) and
area (about 56\%) are occupied by the propulsion system, where each
fan requires an unobstructed duct from side to side.

The most prominent feature of these spaces is the negligible effect of
gravity, called microgravity, resulting from the orbital motion of the
space station around the planet.
We will exploit this feature in
the following way. Most multirotor orient their propellers
vertically along parallel axes. This maximizes thrust vertically in
order to compensate for the gravity force. However, in doing so they
need to tilt in order to move sideways. Moreover, most of the energy
is spent in compensating for gravity. In its absence, not
only there is no need to compensate for gravity, but also the
energy consumption is expected to be significantly less. We propose
to orient the propeller rotation axes non-parallel among them, in
such a way we obtain a fully holonomic vehicle. We expect to
increase maneuverability, being a particularly relevant feature inside
confined spaces such as the interior of a space station.

Holonomic multirotors have been proposed in the past, but the
literature is scarce. In~\cite{jiang13,voyles12} a dexterous
hexarotor has been proposed, where the holonomic kinematics is
used for dexterous manipulation. However, this vehicle is designed for
Earth applications, and thus a trade-off is necessary between dexterity
and gravity compensation. Our design is also based in an hexarotor,
but since we consider microgravity, we will design the propeller
orientation in order to maximize maneuverability along all
directions. Related work in design optimization of hexrotors can be
found in~\cite{rajappa15,kiso15} while in~\cite{langkamp11} a variable
pitch solution is proposed.

The goal of this paper is to present a design of a aerial robot for
microgravity environments. We target a modular and multipurpose
vehicle, providing, for instance, (1)~telepresence, requiring the
vehicle to manoeuvre in such a way the screen and camera is pointed
towards the party onboard, and (2)~mobile manipulation, requiring
force closure in order to compensate the reaction force at the end
effector. Our design also includes a docking port for docking into a
charging station and stacking of multiple units for, e.g., combined
thrust.

This paper is structured as follows: the proposed design is detailed in
Section~\ref{sec:vehicle-design}, in Section~\ref{sec:propulsion} we
derive a parametric model of the propulsion system, which is then used
in Section~\ref{sec:design-parameters} for the optimization its
parameters. Section~\ref{sec:posit-attit-contr} presents a dynamical
model of the vehicle together with a
convergent controller for this holonomic vehicle, and a
preliminary validation in a realistic simulation environment is
presented in Section~\ref{sec:prel-valid}, while
Section~\ref{sec:concl-future-work} wraps up the paper with concluding
remarks and a reference for future work.

This paper builds upon and extends an early draft available in
arXiv~\cite{yoda:arxiv16}.

%%% Local Variables: 
%%% mode: latex
%%% TeX-master: "main"
%%% End: 
